\documentclass[10pt,a4paper]{amsart}
\usepackage[margin=19mm]{geometry}
\usepackage{amsmath,amssymb,mathtools}
\usepackage[scr=rsfs,cal=euler]{mathalfa}
\usepackage{xfrac}
\usepackage[unicode,hidelinks,bookmarks=false]{hyperref}
\newcommand{\ie}{i.e.\ }
\newcommand{\cf}{cf.\ }
\newcommand{\resp}{respectively\ }
\newcommand{\Subsection}{\subsection}
\providecommand{\nobreakdash}{\nobreak\hbox{-}}
\setlength{\emergencystretch}{2em}
\multlinegap0pt
\usepackage{bm}
\usepackage{bbm}
\usepackage{dsfont}
\usepackage{xspace}
\usepackage{cancel}
\usepackage{mathtools}
\usepackage{amsthm}
\makeatletter
\@ifundefined{Lemma}{\newtheorem{Lemma}{Lemma}}{}
\@ifundefined{Remark}{\newtheorem{Remark}{Remark}}{}
\makeatother
\makeatletter
\expandafter\ifx\csname Proof\endcsname\relax
\newenvironment {Proof} {\noindent {\bf Proof.}}{\hspace*{\fill}$\Box$\par\vspace{4mm}}
\fi
\makeatother
\title[The average distance of spanning trees in terms of independe]{The average distance of spanning trees in terms of independence number}
\author{Zhibin Du, Xuli Qi}
\date{Source: arXiv:2605.04725v1 (CC BY 4.0)}
\begin{document}
\maketitle

\noindent\emph{Source and licence.} The average distance of spanning trees in terms of independence number. Version arXiv:2605.04725v1 (CC BY 4.0), distributed under the Creative Commons Attribution 4.0 International licence, as recorded in the arXiv licence metadata for this exact version and shown on the abstract page https://arxiv.org/abs/2605.04725v1. Full rights notice: licenses/arxiv-2605.04725-CC-BY-4.0.txt.\par\smallskip

\noindent\textit{Editorial scope.} Counted unit: the Lemma whose source label is \texttt{le-key} in the article (source0.tex lines 668--674, source label \texttt{le-key}), delivered together with the article's own complete original proof of it (source0.tex lines 676--754, one \texttt{proof} environment of 79 lines). No mathematics is written, repaired or re-derived: statement and proof are the article's own text line for line. Required context retained uncounted: le-final-add (lines 343--355); le-add (lines 476--482); diameter (lines 532--537); remark-final-add (lines 554--559); le-final-2 (lines 562--567); le-extra (lines 616--621). Every definition, display and label of those lines is retained unchanged. Omitted from this package and not redistributed: 1--342, 356--475, 483--531, 538--553, 560--561, 568--615, 622--667, 675--675, 755--996. Nothing inside a retained excerpt is omitted or altered: every retained body is delivered exactly as the article's lines are written, so no omission receipt of any kind accompanies this package. Citations: the retained excerpts carry no citation command at all, so no bibliography entry of the article is transcribed and no reference is redistributed. Reference adaptation: none was needed and none was made; every reference inside the delivered unit resolves inside it, so no unresolved reference or citation is produced. Printed slips kept literal: no printed word, symbol, hypothesis or number of the counted statement or of its proof was altered. Layout and class: the article's own class is \texttt{article} with packages including amsmath, mathrsfs, amssymb, amsfonts, epsfig, indentfirst; the wrapper is the lane's pinned \texttt{amsart} editorial wrapper at 10pt on A4, which restates the theorem-like environments the retained text needs and adds the article's own macro definitions that the retained text needs (none), with extra wrapper packages (bm, bbm, dsfont, xspace, cancel, mathtools). The delivered numbering is the unit's own. Assets: no retained span contains a figure environment, a graphics-inclusion command, a PDF-page inclusion, a table or a TikZ picture; no image file is copied into this package, the package declares no asset dependency and no image file is redistributed with it. Licence: version arXiv:2605.04725v1 of this article is distributed under the Creative Commons Attribution 4.0 International licence, as recorded in the arXiv licence metadata for this exact version and shown on the abstract page https://arxiv.org/abs/2605.04725v1. A full rights notice is delivered at licenses/arxiv-2605.04725-CC-BY-4.0.txt. Characters: every retained excerpt is pure ASCII, so the delivered engine, XeLaTeX with its default Latin Modern text font, reports no missing character.
\begin{Lemma} \label{le-final-add}
Let $G$ be a connected graph of order $n$ and independence number $\alpha$, where $n \ge 2 \alpha$. Then we can find a spanning tree $T$ of $G$ such that $T \in \mathcal{T}_{t, n-t}$, for $t \le 2 \alpha - 1$. In particular, when $t = 2 \alpha - 1$, say $T$ is formed from $\bar{T}$ of order $2 \alpha - 1$ by attaching $n - 2 \alpha + 1$ pendent vertices to some vertices of $\bar{T}$,
\begin{itemize}

\item[(i)]
when $\bar{T}$ is not a path, every pendent path or internal path of $\bar{T}$ is of even length;

\item[(ii)]

when $\bar{T}$ is a path, say $\bar{T} = v_1 v_2 \dots v_{2 \alpha - 1}$, the $n - 2 \alpha + 1$ pendent vertices can only be attached to the vertices among $v_1, v_3, \dots, v_{2 \alpha - 1}$.

\end{itemize}
\end{Lemma}

\begin{Lemma} \label{le-add}
Let $T \in \mathcal{T}_{t,n-t}$, where $t \le 2 k- 2$, $2 \le t \le n - 2$, and $k \ge 2$.
Then
$$
\mu (T) < k + \frac{1}{2}.
$$
\end{Lemma}

\begin{Lemma} \label{diameter}
Let $T$ be a tree obtained from the path $P_{2k-1}$ by attaching $b \ge 1$ pendent vertices to some vertices of $P_{2k-1}$, where $k \ge 2$. If the diameter of $T$ is $2k-2$ or $2k-1$, then
$$
\mu (T) < k + \frac{1}{2}.
$$
\end{Lemma}

\begin{Remark} \label{remark-final-add}
In fact, the above proof indicates that when the number of pendent vertices of $T$ is $n - 2k + 2$,
$$
\mu (T) < k + \frac{1}{2}.
$$
\end{Remark}

\begin{Lemma} \label{le-final-2}
Let $T$ be a tree obtained from the path $P_{2k} = v_1 v_2 \dots v_{2k}$ by attaching $n-2k$ pendent vertices to some vertices of $P_{2k}$, except $v_2$, where $k \ge 2$. If there are at most $\frac{n-2k+1}{k}$ pendent vertices attached to $v_1$, then
$$
\mu (T) < k + \frac{1}{2} + \frac{4(k-1)}{k^2}.
$$
\end{Lemma}

\begin{Lemma} \label{le-extra}
Assume that $T$ is obtained from a tree $H_T$ of order $2k-1$ by attaching $n-2k+1$ pendent vertices to some vertices of $H_T$, i.e., $T \in \mathcal{T}_{2k-1,n-2k+1}$, where $n \ge 2k$ and $k \ge 4$. If $H_T$ is not a path, and each pendent path or internal path in $H_T$ is of even length,  then
$$
\mu (T) < k + \frac{1}{2} + \frac{1}{2(2k-5)}.
$$
\end{Lemma}

\begin{Lemma} \label{le-key}
Let $G$ be a connected graph of order $n$ and independence number $\alpha$, where $n \ge 2 \alpha$ and $\alpha \ge 2$. Assume that $T$ is a spanning tree of $G$ obtained from the path $P_{2\alpha-1} = v_1 v_2 \dots v_{2\alpha-1}$ by attaching $n - 2\alpha + 1$ pendent vertices to some vertices among $v_1, v_3, \dots, v_{2\alpha-1}$.
Then there exists a spanning tree $T^*$ of $G$ such that
$$
\mu(T^*) < \alpha + \frac{1}{2} + \frac{4(\alpha-1)}{\alpha^2}.
$$
\end{Lemma}

\begin{Proof}
Clearly, $T \in \mathcal{T}_{2\alpha-1,n - 2\alpha + 1}$.
On the other hand, note that the diameter of $T$ is possibly $2\alpha - 2$, $2\alpha - 1$, or $2\alpha$. If its diameter is $2\alpha - 2$ or $2\alpha - 1$, then
from Lemma \ref{diameter},
$$
\mu (T) < \alpha + \frac{1}{2}.
$$
So we may assume that the diameter of $T$ is $2\alpha$, it implies that there are some pendent vertices attached to $v_1$ and $v_{2\alpha-1}$ in $T$. Clearly, $T$ has $n - 2\alpha + 1$ pendent vertices.

\noindent {\bf Case 1.} There is some edge joining $v_i$ and $v_j$ in $G$, where $|j - i| \ge 2$.


From Lemma \ref{le-final-add} (ii) and its proof, we know that the vertices $v_1, v_3, \dots, v_{2\alpha - 1}$ form an independent set of $G$. Thus we can assume that either $i$ and $j$ are both even, or $i$ is odd and $j$ is even.

First suppose that $d_T(v_{i+1}) = 2$. Consider $T_1 = T - v_i v_{i+1} + v_i v_j$. Clearly, $T_1$ is still a spanning tree of $G$, and $T_1$ contains one more pendent vertex than $T$, i.e., there are $n - 2\alpha + 2$ pendent vertices in $T_1$. Now from Remark \ref{remark-final-add}, we get
$$
\mu (T_1) < \alpha + \frac{1}{2}.
$$

Similarly, when $d_T(v_{j-1}) = 2$, consider $T_2 = T - v_j v_{j-1} + v_j v_i$ this time, as in last paragraph, we have
$$
\mu (T_2) < \alpha + \frac{1}{2}.
$$

Now there is still one remaining case, i.e., $d_T(v_{i+1}), d_T(v_{j-1}) \ge 3$, i.e., there are some pendent vertices attached to $v_{i+1}$ and $v_{j-1}$ in $T$. It implies that $i$ and $j$ are both even, according to the hypothesis of the structure of $T$. Consider again $T_1 = T - v_i v_{i+1} + v_i v_j$. Notice that if we delete all the pendent vertices of $T_1$, it would lead to a tree, say $T'$, of order $2\alpha-1$ obtained by identifying one end vertices of three paths $P_{k_1}$, $P_{k_2}$, $P_{k_3}$, where $k_1 + k_2 + k_3 = 2\alpha + 1$, $k_1, k_2, k_3 \ge 2$.



On one hand, if at least one of $k_1, k_2, k_3$ is even, then $T'$ has at least one pendent path of odd length, from Lemma \ref{le-final-add} (i), it infers that $T_1 \in \mathcal{T}_{t,n - t}$ with $t \le 2\alpha-2$. Now together with Lemma \ref{le-add},
we have
$$
\mu (T_1) < \alpha + \frac{1}{2}.
$$
On the other hand, if each of $k_1, k_2, k_3$ is odd, i.e., all the pendent paths of $T'$ are of even lengths, then by Lemma \ref{le-extra},
$$
\mu (T_1) < \alpha + \frac{1}{2} + \frac{1}{2(2\alpha-5)}.
$$



\noindent {\bf Case 2.} The induced subgraph of $G$ induced by the vertices $v_1, v_2, \dots, v_{2\alpha-1}$ is still $P_{2\alpha-1}$.

Note that we may assume that there is at least one pendent neighbor of $v_1$ in $T$, say $u$, which is adjacent to none of $v_2, v_3, \dots, v_{2\alpha-1}$ in $G$. Otherwise,
every neighbor of $v_1$ in $T$ is adjacent to at least one vertex of $v_2, v_3, \dots, v_{2\alpha-1}$,
we can always construct another spanning tree $T_3$, which is obtained from $P_{2\alpha-1}- v_1 = v_2 v_3 \dots v_{2\alpha-1}$ (of order $2\alpha-2$) by attaching $n - 2\alpha + 2$ pendent vertices to some vertices of $ v_2, v_3, \dots, v_{2\alpha-1}$, i.e., $T_3 \in \mathcal{T}_{2\alpha-2,n - 2\alpha + 2}$. Then by Lemma \ref{le-add},
$$
\mu (T_3) < \alpha + \frac{1}{2}.
$$

Similarly, we may also assume that there is at least one pendent neighbor of $v_{2\alpha-1}$ in $T$, say $w$, which is adjacent to none of $v_1, v_2, \dots, v_{2\alpha-2}$ in $G$.

If $u$ and $w$ are not adjacent in $G$, then $\{u, w, v_2, v_4, \dots, v_{2\alpha-2}\}$ would form an independent set of $G$, however, it contains $\alpha + 1$ vertices, which is a contradiction. Thus $u$ and $w$ are adjacent in $G$.

Denote by $a_k$ the number of pendent vertices attached to $v_{2k-1}$ in $T$, where $1 \le k \le \alpha$. Clearly,
$$
a_1 + a_2 + \dots + a_{\alpha} = n - 2 \alpha + 1.
$$
Note that we can always find some $t$, where $2 \le t \le \alpha$, such that $a_t \le \frac{n - 2 \alpha + 1}{\alpha}$, just need to assume that
$$
a_t = \min \{ a_1, a_2, \dots, a_{\alpha} \}.
$$
Here we exclude the case $t = 1$ (which is equivalent to $t = \alpha$) for convenience.
Let $T_4 = T - v_{2t - 2} v_{2t - 1} + u w$. Clearly, $T_4$ is still a spanning tree of $G$. More precisely, $T_4$ is a tree obtained from the path
$$
v_{2t-1} v_{2t} \dots v_{2\alpha - 1} w u v_1 v_2 \dots v_{2t - 3}
$$
on $2\alpha$ vertices by attaching some pendent vertices to this path. In particular, no pendent vertex outside such path on $2\alpha$ vertices is attached to $v_{2t}$ when $2 \le t \le \alpha -1$ ($w$ when $t = \alpha$, respectively) in $T_4$ (from the hypothesis of the structure of $T$). Now by Lemma \ref{le-final-2},
$$
\mu (T_4) < \alpha + \frac{1}{2} + \frac{4(\alpha-1)}{\alpha^2}
$$
follows.


Note that
$$
\frac{1}{2(2\alpha-5)} < \frac{4(\alpha-1)}{\alpha^2}.
$$
Combining Cases $1$ and $2$, by setting $T^* = T, T_1, T_2, T_3$ or $T_4$ accordingly, we would get the desired result.
\end{Proof}

\end{document}
